\section{Claim boundary and conclusion}
\label{sec:claim-boundary}

\subsection{What is proved}

The following statements are unconditional for arbitrary complex
sequences:
\begin{enumerate}[label=\textup{(\arabic*)},leftmargin=2.5em]
 \item the exact discrete Abel identity
 \[
  S_N=NL_N-\sum_{m<N}L_m;
 \]
 \item the necessary-and-sufficient signed Ces\`aro criterion
 \[
  S_N/N\to0\iff
  L_N-\frac1N\sum_{m<N}L_m\to0;
 \]
 \item convergence of \(L_N\) is sufficient, with the quantitative
 bounds \eqref{eq:convergence-rate} and
 \eqref{eq:tail-convergence-rate};
 \item terminal-half logarithmic cancellation and vanishing dyadic
 logarithmic variation are sufficient, with the explicit bounds
 \eqref{eq:geometric-local-global} and
 \eqref{eq:variation-convolution};
 \item a fixed-power terminal-half bound below exponent one
 transfers without a fixed-power loss;
 \item the bounded dyadic sequence
 \eqref{eq:dyadic-sequence} has normalized logarithmic cancellation
 on every growing multiplicative window but no ordinary
 cancellation.
\end{enumerate}
These are \(\Lzero\) identities, inequalities, and a
structure-free \(\Lzero\) countermodel.

\Cref{thm:literal-certificate} is also an exact finite implication.
Its status becomes \(\Lone\) only after
\eqref{eq:literal-reassembly} is verified for the complete raw hard
physical packet.

\subsection{What is not proved}

This paper does not prove:
\begin{enumerate}[label=\textup{(\alph*)},leftmargin=2.5em]
 \item terminal-half \(o(1)\) or fixed-power cancellation for the
 actual growing affine M\"obius family;
 \item uniformity of Tao's fixed-form theorem in the polynomially
 growing physical coefficients;
 \item selection of the prescribed \(h_0\) from an averaged-shift
 theorem;
 \item removal of the exceptional scales or interval origins in an
 almost-all theorem;
 \item cancellation on the \(O(1)\)-length physical fibers;
 \item a complete physical outer-reassembly bound;
 \item a determinant-energy lower bound for \(D_X\);
 \item a fixed-\(h_0\) Hardy--Littlewood asymptotic;
 \item a breach of the general sieve parity barrier;
 \item a prime-pair lower bound or the twin-prime conjecture.
\end{enumerate}
No conclusion is specialized from a general fixed nonzero \(h_0\)
to \(h_0=2\).

\subsection{Conclusion}

The logarithmic-to-ordinary gap is not an informal change of weight.
For the same ordered coefficient sequence it is measured exactly by
\[
 \cT_N(L)=
 L_N-\frac1N\sum_{m<N}L_m.
\]
Global boundedness of \(L_N\), or even normalized cancellation on
every terminal window whose ratio tends to infinity, does not force
this defect to vanish.  The smallest exact target is therefore the
signed defect itself.  Terminal-half cancellation and dyadic
variation are stronger, quantitative certificates that can be
checked without ambiguity.

For the TPC route this produces a clean fork.  A positive advance
must prove one of these certificates uniformly for the literal
growing fixed-\(h_0\) affine fibers and survive the full outer
reassembly.  A negative advance may prove that a proposed
logarithmic input cannot meet that certificate under its native
quantifiers.  Either outcome must retain the actual coefficient,
global normalization, and the strict endpoint rule
\(\Lambda_{\mathrm{phys}}<1/400\).  The present paper identifies that
gate; it does not cross the unresolved arithmetic parity boundary.
